\courseunit{1}
\emergencystretch=2em

\begin{problem}{1}[\textbf{Basic properties of growth rates.}]
Use the fact that the growth rate of a variable equals the time derivative of its log to show:
\begin{enumerate}[label=(\alph*)]
    \item The growth rate of the product of two variables equals the sum of their growth rates.
    That is, if \(Z\left(t\right)=X\left(t\right)Y\left(t\right)\), then
    \begin{align*}
        \dfrac{\dot Z\left(t\right)}{Z\left(t\right)} = 
        \dfrac{\dot X\left(t\right)}{X\left(t\right)} + 
        \dfrac{\dot Y\left(t\right)}{Y\left(t\right)}
    \end{align*}
    \item The growth rate of the ratio of two variables equals the difference of their growth rates.
    That is, if \(Z\left(t\right)=X\left(t\right)/Y\left(t\right)\), then
    \begin{align*}
        \dfrac{\dot Z\left(t\right)}{Z\left(t\right)} = 
        \dfrac{\dot X\left(t\right)}{X\left(t\right)} - 
        \dfrac{\dot Y\left(t\right)}{Y\left(t\right)}
    \end{align*}
    \item If \(Z\left(t\right)=a\,X\left(t\right)^{\alpha}\), then
    \(\dot{Z}\left(t\right)/Z\left(t\right)=\alpha\,\dot{X}\left(t\right)/X\left(t\right)\).
\end{enumerate}
\end{problem}

\begin{problem}{2}
Suppose that the growth rate of some variable, \(X\), is constant and equal to \(a>0\) from time 0
to time \(t_1\); drops to 0 at time \(t_1\); rises gradually from 0 to \(a\) from time \(t_1\)
to time \(t_2\); and is constant and equal to \(a\) after time \(t_2\).
\begin{enumerate}[label=(\alph*)]
    \item Sketch a graph of the growth rate of \(X\) as a function of time.
    \item Sketch a graph of \(\ln X\) as a function of time.
\end{enumerate}
\end{problem}

\begin{problem}{3}
Describe how, if at all, each of the following developments affects the break-even and actual
investment lines in our basic diagram for the Solow model:
\begin{enumerate}[label=(\alph*)]
    \item The rate of depreciation falls.
    \item The rate of technological progress rises.
    \item The production function is Cobb--Douglas, \(f\left(k\right)=k^{\alpha}\), and
    capital's share, \(\alpha\), rises.
    \item Workers exert more effort, so that output per unit of effective labor for a given value
    of capital per unit of effective labor is higher than before.
\end{enumerate}
\end{problem}

\begin{problem}{4}
Consider an economy with technological progress but without population growth that is on its balanced
growth path. Now suppose there is a one-time jump in the number of workers.
\begin{enumerate}[label=(\alph*)]
    \item At the time of the jump, does output per unit of effective labor rise, fall, or stay
    the same? Why?
    \item After the initial change (if any) in output per unit of effective labor when the new
    workers appear, is there any further change in output per unit of effective labor? If so,
    does it rise or fall? Why?
    \item Once the economy has again reached a balanced growth path, is output per unit of
    effective labor higher, lower, or the same as it was before the new workers appeared? Why?
\end{enumerate}
\end{problem}

\begin{problem}{5}
Suppose that the production function is Cobb--Douglas.
\begin{enumerate}[label=(\alph*)]
    \item Find expressions for \(k^*\), \(y^*\), and \(c^*\) as functions of the parameters
    of the model, \(s\), \(n\), \(\delta\), \(g\), and \(\alpha\).
    \item What is the golden-rule value of \(k\)?
    \item What saving rate is needed to yield the golden-rule capital stock?
\end{enumerate}
\end{problem}

\begin{problem}{6}
Consider a Solow economy that is on its balanced growth path. Assume for simplicity that there is
no technological progress. Now suppose that the rate of population growth falls.
\begin{enumerate}[label=(\alph*)]
    \item What happens to the balanced-growth-path values of capital per worker, output per worker,
    and consumption per worker? Sketch the paths of these variables as the economy moves to its new
    balanced growth path.
    \item Describe the effect of the fall in population growth on the path of output (that is,
    total output, not output per worker).
\end{enumerate}
\end{problem}

\begin{problem}{7}
Find the elasticity of output per unit of effective labor on the balanced growth path, \(y^*\),
with respect to the rate of population growth, \(n\). If \(\alpha_K\left(k^*\right)=1/3\),
\(g=2\%\), and \(\delta=3\%\), by about how much does a fall in \(n\) from 2 percent to
1 percent raise \(y^*\)?
\end{problem}

\begin{problem}{8}
Suppose that investment as a fraction of output in the United States rises permanently from 0.15
to 0.18. Assume that capital's share is \(1/3\).
\begin{enumerate}[label=(\alph*)]
    \item By about how much does output eventually rise relative to what it would have been without the rise in investment?
    \item By about how much does consumption rise relative to what it would have been without the rise in investment?
    \item What is the immediate effect of the rise in investment on consumption? About how long
    does it take for consumption to return to what it would have been without the rise in investment?
\end{enumerate}
\end{problem}

\begin{problem}{9}[\textbf{Factor payments in the Solow model.}]
Assume that both labor and capital are paid their marginal products. Let \(w\) denote
\(\partial F\left(K,AL\right)/\partial L\) and \(r\) denote
\(\left[\partial F\left(K,AL\right)/\partial K\right]-\delta\).
\begin{enumerate}[label=(\alph*)]
    \item Show that the marginal product of labor, \(w\), is \(A\left[f\left(k\right)-k\,f'\left(k\right)\right]\).
    \item Show that if both capital and labor are paid their marginal products, constant returns
    to scale imply that the total amount paid to the factors of production equals total net output.
    That is, show that under constant returns, \(wL+rK=F\left(K,AL\right)-\delta K\).
    \item The return to capital (\(r\)) is roughly constant over time, as are the shares of output
    going to capital and to labor. Does a Solow economy on a balanced growth path exhibit these
    properties? What are the growth rates of \(w\) and \(r\) on a balanced growth path?
    \item Suppose the economy begins with a level of \(k\) less than \(k^*\). As \(k\) moves
    toward \(k^*\), is \(w\) growing at a rate greater than, less than, or equal to its growth
    rate on the balanced growth path? What about \(r\)?
\end{enumerate}
\end{problem}

\begin{problem}{10}
\problemsource{\parencite{piketty2014capital,piketty-zucman2014capital,rognlie2015deciphering}}
This problem asks you to use a Solow-style model to investigate some ideas that have been discussed
in the context of Thomas Piketty's recent work. Consider an economy described by the assumptions of the Solow model, except that
factors are paid their marginal products (as in the previous problem), and all labor income is
consumed and all other income is saved. Thus,
\(C\left(t\right)=L\left(t\right)\left[\partial Y\left(t\right)/\partial L\left(t\right)\right]\).
\begin{enumerate}[label=(\alph*)]
    \item Show that the properties of the production function and our assumptions about the behavior
    of \(L\) and \(A\) imply that the capital-output ratio, \(K/Y\), is rising if and only if
    the growth rate of \(K\) is greater than \(n+g\)--that is, if and only if \(k\) is rising.
    \item Assume that the initial conditions are such that \(\partial Y/\partial K\) at \(t=0\)
    is strictly greater than \(n+g+\delta\). Describe the qualitative behavior of the capital-output
    ratio over time. (For example, does it grow or fall without bound? Gradually approach some
    constant level from above or below? Something else?) Explain your reasoning.
    \item Many popular summaries of Piketty's work describe his thesis as: Since the return to
    capital exceeds the growth rate of the economy, the capital-output ratio tends to grow without
    bound. By the assumptions in part (b), this economy starts in a situation where the return to
    capital exceeds the economy's growth rate. If you found in (b) that \(K/Y\) grows without
    bound, explain intuitively whether the driving force of this unbounded growth is that the return
    to capital exceeds the economy's growth rate. Alternatively, if you found in (b) that \(K/Y\)
    does not grow without bound, explain intuitively what is wrong with the statement that the return
    to capital exceeding the economy's growth rate tends to cause \(K/Y\) to grow without bound.
    \item Suppose \(F\left(\,\cdot\,\right)\) is Cobb--Douglas and that the initial situation is as in
    part (b). Describe the qualitative behavior over time of the share of net capital income (that is,
    \(K\left(t\right)\left[\partial Y\left(t\right)/\partial K\left(t\right)-\delta\right]\))
    in net output (that is, \(Y\left(t\right)-\delta K\left(t\right)\)). Explain your reasoning.
    Is the common statement that an excess of the return to capital over the economy's growth rate
    causes capital's share to rise over time correct in this case?
\end{enumerate}
\end{problem}

\begin{problem}{11}
Consider Problem 1.10. Suppose there is a marginal increase in \(K\).
\begin{enumerate}[label=(\alph*)]
    \item Derive an expression (in terms of \(K/Y\), \(\delta\), the marginal product of capital
    \(F_K\), and the elasticity of substitution between capital and effective labor in the gross
    production function \(F\left(\,\cdot\,\right)\)) that determines whether a marginal increase
    in \(K\) increases, reduces, or has no effect on the share of net capital income in net output.
    \item Suppose the capital-output ratio is 3, \(\delta=3\%\), and the rate of return on capital
    (\(F_K-\delta\)) = 5\%. How large must the elasticity of substitution in the gross production
    function be for the share of net capital income in net output to rise when \(K\) rises?
\end{enumerate}
\end{problem}

\begin{problem}{12}
Consider the same setup as at the start of Problem 1.10: the economy is described by the assumptions
of the Solow model, except that factors are paid their marginal products and all labor income is
consumed and all other income is saved. Show that the economy converges to a balanced growth path,
and that the balanced-growth-path level of \(k\) equals the golden-rule level of \(k\). What is
the intuition for this result?
\end{problem}

\begin{problem}{13}
Go through steps analogous to those in equations (1.29)--(1.32) to find how quickly \(y\) converges
to \(y^*\) in the vicinity of the balanced growth path.
\par \remarkaccent{Hint:} Since \(y=f\left(k\right)\), we can write \(k=g\left(y\right)\), where \(g\left(\,\cdot\,\right)=f^{-1}\left(\,\cdot\,\right)\).
\end{problem}

\begin{problem}{14}[\textbf{Embodied technological progress.}]
\problemsource{\parencite{solow1960investment,sato1966adjustment}}
One view of technological progress is that the
productivity of capital goods built at \(t\) depends on the state of technology at \(t\) and is
unaffected by subsequent technological progress. This is known as \emph{embodied technological
progress} (technological progress must be ``embodied'' in new capital before it can raise output).
This problem asks you to investigate its effects.
\begin{enumerate}[label=(\alph*)]
    \item As a preliminary, let us modify the basic Solow model to make technological progress
    capital-augmenting rather than labor-augmenting. So that a balanced growth path exists, assume
    that the production function is Cobb--Douglas:
    \begin{align*}
    Y\left(t\right)
    &=\left[A\left(t\right)K\left(t\right)\right]^{\alpha}L\left(t\right)^{1-\alpha}
    \end{align*}
    Assume that \(A\) grows at rate \(\mu\): \(\dot{A}\left(t\right)=\mu A\left(t\right)\).
    Show that the economy converges to a balanced growth path, and find the growth rates of \(Y\)
    and \(K\) on the balanced growth path.
    \par \remarkaccent{Hint:} Show that we can write
    \(Y/\left(A^{\phi}L\right)\) as a function of \(K/\left(A^{\phi}L\right)\), where
    \(\phi=\alpha/\left(1-\alpha\right)\). Then analyze the dynamics of
    \(K/\left(A^{\phi}L\right)\).
    \item Now consider embodied technological progress. Specifically, let the production function be
    \(Y\left(t\right)=J\left(t\right)^{\alpha}L\left(t\right)^{1-\alpha}\), where \(J\left(t\right)\)
    is the effective capital stock. The dynamics of \(J\left(t\right)\) are given by
    \(\dot{J}\left(t\right)=s\,A\left(t\right)Y\left(t\right)-\delta J\left(t\right)\).
    The presence of the \(A\left(t\right)\) term in this expression means that the productivity
    of investment at \(t\) depends on the technology at \(t\). Show that the economy converges
    to a balanced growth path. What are the growth rates of \(Y\) and \(J\) on the balanced
    growth path?
    \par \remarkaccent{Hint:} Let \(\bar{J}\left(t\right)\equiv J\left(t\right)/A\left(t\right)\). Then use the same approach as in (a), focusing on \(\bar{J}/\left(A^{\phi}L\right)\) instead of \(K/\left(A^{\phi}\,L\right)\).
    \item What is the elasticity of output on the balanced growth path with respect to \(s\)?
    \item In the vicinity of the balanced growth path, how rapidly does the economy converge to the balanced growth path?
    \item Compare your results for (c) and (d) with the corresponding results in the text for
    the basic Solow model.
\end{enumerate}
\end{problem}

\begin{problem}{15}
Consider a Solow economy on its balanced growth path. Suppose the growth-accounting techniques
described in Section 1.7 are applied to this economy.
\begin{enumerate}[label=(\alph*)]
    \item What fraction of growth in output per worker does growth accounting attribute to growth
    in capital per worker? What fraction does it attribute to technological progress?
    \item How can you reconcile your results in (a) with the fact that the Solow model implies that
    the growth rate of output per worker on the balanced growth path is determined solely by the rate
    of technological progress?
\end{enumerate}
\end{problem}

\begin{problem}{16}
\begin{enumerate}[label=(\alph*)]
    \item In the model of convergence and measurement error in equations (1.39) and (1.40), suppose
    the true value of \(b\) is \(-1\). Does a regression of
    \(\ln\left(Y/N\right)_{1979}-\ln\left(Y/N\right)_{1870}\) on a constant and
    \(\ln\left(Y/N\right)_{1870}\) yield a biased estimate of \(b\)? Explain.
    \item Suppose there is measurement error in measured 1979 income per capita but not in 1870
    income per capita. Does a regression of
    \(\ln\left(Y/N\right)_{1979}-\ln\left(Y/N\right)_{1870}\) on a constant and
    \(\ln\left(Y/N\right)_{1870}\) yield a biased estimate of \(b\)? Explain.
\end{enumerate}
\end{problem}

\begin{problem}{17}
Derive equation (1.50).
\par \remarkaccent{Hint:} Follow steps analogous to those in equations [1.47] and [1.48].
\end{problem}
